Em \enquote{As heterotopias}, Foucault começa por evocar planetas que não se encontram em mapa algum e, pouco adiante, encontra um oceano na cama dos pais, onde as crianças nadam entre as cobertas (pp. 19–20). A cama não deixa de ser cama para tornar-se oceano: leio nessa brincadeira o que Foucault descreve adiante como a possibilidade de reunir, num mesmo lugar real, espaços que habitualmente seriam incompatíveis (p. 24). Chamo Cecinepá de heterotopia especulativa porque quero experimentar com esta carta algo semelhante: as ilhas são imaginadas, mas é aqui, na leitura, que o nariz do perfumista e o cromatograma do químico podem se encontrar e suscitar uma pergunta que leve ambos de volta às amostras.